\par\Needspace{6\baselineskip}\medskip
\SourceLeaf{136}\pdfbookmark[2]{98.}{article-098}
\begin{center}98.\end{center}
Die in den vorangehenden Artikeln entwickelte Lösung unserer Aufgabe für die Ellipse ließe sich auch auf die Parabel und die Hyperbel übertragen. Dazu würde man die Parabel als Ellipse auffassen, in der $a$ und $b$ unendliche Größen wären und $\varphi=90^0$ sowie schließlich $E$, $E'$, $g$ und $G=0$ gälten. Ebenso würde man die Hyperbel als Ellipse auffassen, in der $a$ negativ und $b$, $E$, $E'$, $g$, $G$ und $\varphi$ imaginär wären. Wir ziehen es jedoch vor, diese Annahmen\raisebox{.45ex}{.} beiseitezulassen und die Aufgabe für jede der beiden Arten von Kegelschnitten gesondert zu behandeln. Die auffallende Analogie zwischen allen drei Arten wird so von selbst sichtbar.

Für die PARABEL behalten wir die Zeichen $p$, $v$, $v'$, $F$, $f$, $r$, $r'$ und $t$ in ihrer bisherigen Bedeutung bei. Aus der Theorie der parabolischen Bewegung folgt
\[
\begin{aligned}
\mathrm{1)}\qquad&\BareRoot{\dfrac p{2r}}=\cos\tfrac12(F-f)\\[.7em]
\mathrm{2)}\qquad&\BareRoot{\dfrac p{2r'}}=\cos\tfrac12(F+f)
\end{aligned}
\]
\[
\begin{aligned}
\frac{2kt}{p^{\frac32}}&=\operatorname{tang}\tfrac12(F+f)-\operatorname{tang}\tfrac12(F-f)\\
&\quad+\tfrac13\operatorname{tang}\tfrac12(F+f)^3-\tfrac13\operatorname{tang}\tfrac12(F-f)^3\\[.5em]
&=\{\operatorname{tang}\tfrac12(F+f)-\operatorname{tang}\tfrac12(F-f)\}\\
&\quad\Bigl\{1+\operatorname{tang}\tfrac12(F+f)\operatorname{tang}\tfrac12(F-f)\\
&\qquad+\tfrac13(\operatorname{tang}\tfrac12(F+f)-\operatorname{tang}\tfrac12(F-f))^2\Bigr\}\\[.5em]
&=\frac{2\sin f\,\mathpunct{.}\BareRoot{rr'}}p
\left\{\frac{2\cos f\,\mathpunct{.}\BareRoot{rr'}}p+\frac{4\sin f^2rr'}{3pp}\right\},
\end{aligned}
\]
\noindent woraus folgt
\[
\mathrm{3)}\qquad kt=\frac{2\sin f\cos f\,\mathpunct{.}\,rr'}{\BareRoot{p}}+\frac{4\sin f^3(rr')^{\frac32}}{3p^{\frac32}}\mathpunct{\cdot}
\]
\noindent Durch Multiplikation der Gleichungen 1 und 2 erhält man ferner
\[
\mathrm{4)}\qquad \frac p{\BareRoot{rr'}}=\cos F+\cos f,
\]
\noindent und durch Addition ihrer Quadrate
\[
\mathrm{5)}\qquad \frac{p(r+r')}{2rr'}=1+\cos F\cos f.
\]
\noindent Eliminiert man daraus $\cos F$, so erhält man
\[
\mathrm{6)}\qquad p=\frac{2rr'\sin f^2}{r+r'-2\cos f\,\mathpunct{.}\BareRoot{rr'}}\mathpunct{\cdot}
\]
\noindent Übernehmen wir also auch hier die Gleichungen 9 und $9^*$ aus Art. 88, die erste für \SourceLeaf{137}positives $\cos f$ und die zweite für negative Werte, so erhalten wir
\[
\begin{aligned}
\mathrm{7)}\qquad&p=\frac{\sin f^2\BareRoot{rr'}}{2l\cos f}\\[.6em]
\mathrm{7^*)}\qquad&p=\frac{\sin f^2\BareRoot{rr'}}{-2L\cos f}\mathpunct{\raisebox{.45ex}{,}}
\end{aligned}
\]
\noindent Setzt man diese Werte in Gleichung 3 ein und behält die Zeichen $m$ und $M$ in der durch die Gleichungen 11 und $11^*$ aus Art. 88 festgelegten Bedeutung bei, so erhält man
\[
\begin{aligned}
\mathrm{8)}\qquad&m=l^{\frac12}+\tfrac43l^{\frac32}\\[.5em]
\mathrm{8^*)}\qquad&M=-L^{\frac12}+\tfrac43L^{\frac32}.
\end{aligned}
\]
\noindent Diese Gleichungen stimmen mit 12 und $12^*$ aus Art. 88 überein, wenn man dort $g=0$ setzt. Nimmt man folglich bei zwei heliozentrischen Orten, denen eine Parabel genügt, eine elliptische Bahn an, so muss sich bei Anwendung der Vorschriften aus Art. 91 sogleich $x=0$ ergeben. Ergibt sich umgekehrt nach diesen Vorschriften $x=\mathpunct{\raisebox{.65ex}{.}}0$, so sieht man leicht, dass sich die angenommene Ellipse als Parabel erweist, da die Gleichungen 1, 16, 17, 19 und 20 die Werte $b=\infty$, $a=\infty$ und $\varphi=90^0$ liefern. Die Bestimmung der Elemente lässt sich dann sehr leicht abschließen. Für $p$ kann man entweder Gleichung 7 dieses Artikels oder Gleichung 18 aus Art. 95 verwenden\footnote{Daraus ist zugleich ersichtlich, dass $y$ und $Y$ in der Parabel dieselben Verhältnisse ausdrücken wie in der Ellipse. Siehe Art. 95.}. Für $F$ ergibt sich aus den Gleichungen 1 und 2 dieses Artikels $\operatorname{tang}\tfrac12 F=\dfrac{\BareRoot{r'}-\BareRoot{r}}{\BareRoot{r'}+\BareRoot{r}}\operatorname{cotang}\tfrac12 f=\sin2\omega\operatorname{cotang}\tfrac12 f$, wenn der Hilfswinkel dieselbe Bedeutung wie in Art. 89 hat.

Bei dieser Gelegenheit bemerken wir noch, dass sich die bekannte Gleichung ergibt, wenn man in Gleichung 3 für $p$ seinen Wert aus Gleichung 6 einsetzt.
\[
kt=\tfrac13(r+r'+\cos f\,\mathpunct{.}\BareRoot{rr'})(r+r'-2\cos f\,\mathpunct{.}\BareRoot{rr'})^{\frac12}\BareRoot{2}.
\]
